\documentclass[10pt]{article}
\usepackage[T1]{fontenc}
\usepackage[utf8]{inputenc}
\usepackage[french]{babel}
\usepackage{amsmath,amssymb,mathrsfs,mathpazo}
\usepackage{geometry}
\geometry{a4paper,margin=2cm}
\usepackage{graphicx,array,tabularx,tabularray,colortbl,multirow,multicol,booktabs,enumitem,pifont,ulem,fancybox,diagbox,cancel}
\usepackage{pstricks,pstricks-add,pst-plot,pst-tree,pst-node,pst-text,pst-func,pst-math,pst-bspline,pst-3dplot,pst-eucl,pst-solides3d}
\let\pscalculate\undefined
\usepackage{pst-bezier}
\usepackage{tikz,pgfplots,tkz-tab}
\usetikzlibrary{arrows,arrows.meta,patterns,calc,positioning,trees,decorations.markings,shapes.geometric,shapes.misc,shapes.arrows,intersections,decorations.pathmorphing,decorations.pathreplacing,backgrounds,shadows,babel}
\pgfplotsset{compat=1.18}
\usepgfplotslibrary{fillbetween}
\usepackage[np]{numprint}
\usepackage{listings,makecell,pas-tableur,twoopt}
\let\barre\undefined
\usepackage[pstricks]{tablvar}
\definecolor{darkgreen}{RGB}{0,100,0}
\providecommand{\red}{\color{red}}
\providecommand{\blue}{\color{blue}}
\providecommand{\green}{\color{green}}
\providecommand{\black}{\color{black}}
\providecommand{\white}{\color{white}}
\providecommand{\gray}{\color{gray}}
\providecommand{\cyan}{\color{cyan}}
\providecommand{\magenta}{\color{magenta}}
\providecommand{\yellow}{\color{yellow}}
\providecommand{\violet}{\color{violet}}
\providecommand{\orange}{\color{orange}}
\providecommand{\euro}{\texteuro}
\providecommand{\EUR}{\texteuro}
\providecommand{\ieme}{\textsuperscript{e}}
\providecommand{\degres}{\ensuremath{^\circ}}
\providecommand{\degree}{\ensuremath{^\circ}}
\providecommand{\vv}[1]{\overrightarrow{#1}}
\providecommand{\Coord}[1]{#1}
\providecommand{\decosix}{}
\providecommand{\starredbullet}{\ensuremath{\bullet}}
\pagestyle{empty}
\setlength{\parindent}{0pt}
\usepackage{adjustbox}\DeclareUnicodeCharacter{2236}{:}\DeclareUnicodeCharacter{2190}{\ensuremath{\leftarrow}}
\providecommand{\R}{\mathbb{R}}
\providecommand{\N}{\mathbb{N}}
\providecommand{\Z}{\mathbb{Z}}
\providecommand{\Q}{\mathbb{Q}}
\providecommand{\D}{\mathbb{D}}
\providecommand{\C}{\mathbb{C}}
\providecommand{\e}{{\mathrm{e}}}
\providecommand{\vect}[1]{\overrightarrow{#1}}
\providecommand{\vectt}[1]{\overrightarrow{\mathrm{#1}}}
\providecommand{\vvt}[1]{\overrightarrow{\mathrm{#1}}}
\providecommand{\np}[2][]{{#2~#1}}
\providecommand{\vv}[1]{{\overrightarrow{#1}}}
\renewcommand{\d}{{\mathrm{d}}}
\providecommand{\degres}{{^\circ}}
\providecommand{\degre}{{^\circ}}
\providecommand{\blue}{\color{blue}}
\providecommand{\red}{\color{red}}
\providecommand{\enskip}{\quad}
\providecommand{\llbracket}{\lbrack\!\lbrack}
\providecommand{\rrbracket}{\rbrack\!\rbrack}
\providecommand{\euro}{{\text{€}}}
\providecommand{\no}{{\text{n°}}}
\providecommand{\interval}[2]{{[#1;#2]}}
\providecommand{\intervFO}[2]{{[#1;#2[}}
\providecommand{\dotfill}{\ldots}

\providecommand{\ds}{\displaystyle}
\providecommand{\textding}[1]{#1}
\providecommand{\Oij}{\(\mathrm{O};\vec{\imath},\vec{\jmath}\)}
\providecommand{\Oijk}{\(\mathrm{O};\vec{\imath},\vec{\jmath},\vec{k}\)}
\providecommand{\Ouv}{\(\mathrm{O};\vec{u},\vec{v}\)}
\providecommand{\intervFF}[2]{\ensuremath{[#1;#2]}}
\renewcommand{\euro}{\ifmmode\text{\texteuro}\else\texteuro\fi}
\DeclareUnicodeCharacter{00B0}{\ifmmode^{\circ}\else\textdegree\fi}
\DeclareUnicodeCharacter{20AC}{\ifmmode\text{\texteuro}\else\texteuro\fi}
% Énoncés : Éducation nationale. Transcription et corrigés : APMEP, Denis Vergès et contributeurs. Découpage et classement : CoopMaths, Sébastien Lozano, Éric Elter et contributeurs. Adaptation de la présentation : Maths Explorer.
% CC BY-SA : https://coopmaths.fr/www/about/
% https://coopmaths.fr/alea/static/bac/2024/tex/bac_2024_06_sujet2_polynesie_3.tex
\begin{document}
\noindent\textit{APMEP / CoopMaths. CC BY-SA. Présentation adaptée par Maths Explorer.}\par\medskip
\medskip

On considère la suite $\left(u_n\right)$ définie par : 

\begin{center} $u_0 = 8$ \:  et pour tout entier naturel \:\:$n,\:\:~u_{n+1} = u_n -\ln\left(\dfrac{u_n}{4}\right)$.\end{center}

\begin{enumerate}
\item 
	\begin{enumerate}
		\item Donner les valeurs arrondies au centième de $u_1$ et $u_2$.
		\item On considère la fonction \texttt{mystere} définie ci-dessous en \textsf{Python}. On admet que, pour tout réel strictement positif \texttt{a}, \texttt{log(a)} renvoie la valeur du logarithme népérien de \texttt{a}.

\begin{center}
\begin{tabular}{|l|}\hline
def mystere(k) :\\
\qquad u = 8\\
\qquad S = 0\\
\qquad for i in range(k) :\\
\qquad \qquad S = S + u\\
\qquad \qquad u = u $-$ log( u / 4 )\\
\qquad return S\\ \hline
\end{tabular}
\end{center}
	
L'exécution de \texttt{mystere(10)} renvoie \texttt{58.44045206721732}. Que représente ce résultat ?
		\item Modifier la fonction précédente afin qu'elle renvoie la moyenne des $k$ premiers termes de la suite $\left(u_n\right)$.
	\end{enumerate}
\item On considère la fonction $f$ définie et dérivable sur $]0~;~+\infty[$ par :

\[ f(x) = x - \ln\left(\dfrac{x}{4}\right).\]

On donne ci-dessous une représentation graphique $\mathcal{C}_f$ de la fonction $f$ pour les valeurs de $x$ comprises entre 0 et 6.
	
\begin{center}
\psset{xunit=2cm,yunit=1.2cm,arrowsize=2pt 3}
\begin{pspicture}(-0.5,-0.25)(6.1,6.1)
\psgrid[gridlabels=0pt,subgriddiv=1,gridwidth=0.3pt]
\psaxes[linewidth=1.25pt,labelFontSize=\scriptstyle]{->}(0,0)(0,0)(6.1,6.1)
\psaxes[linewidth=1.25pt,labelFontSize=\scriptstyle](0,0)(0,0)(6.1,6.1)
\psplot[plotpoints=2000,linewidth=1.25pt,linecolor=red]{0.01}{6}{x  x 4 div ln sub}
\uput[u](5.5,5.2){\red $\mathcal{C}_f$}
\end{pspicture}
\end{center}
	
Étudier les variations de $f$ sur $]0~;~+\infty[$ et dresser son tableau de variations.
	
\emph{On précisera la valeur exacte du minimum de $f$ sur $]0~;~+\infty[$. Les limites ne sont pas demandées.}
\end{enumerate}

Dans la suite de l'exercice, on remarquera que pour tout entier naturel $n$,\: $u_{n+1} = f\left(u_n\right)$.

\begin{enumerate}[resume]
\item 
	\begin{enumerate}
		\item Démontrer, par récurrence, que pour tout entier naturel $n$, on a : 

		\[ 1 \leqslant u_{n+1} \leqslant u_n. \]
		\item En déduire que la suite $\left(u_n\right)$ converge vers une limite réelle.

On note $\ell$ la valeur de cette limite 
		\item Résoudre l'équation $f(x) = x$.
		\item En déduire la valeur de $\ell$.
	\end{enumerate}
\end{enumerate}

\bigskip

\end{document}
